\chapter{现象学读出：记忆、时间感与体验报告的计算边界}
\label{chp:phenomenological_physics}
记忆、时间感与感受质处在计算模型能够解释和主体体验不能被轻率还原的交界处。我们可以测量活动维持、结构学习、检索行为、事件排序和体验报告，却不能仅凭内部状态稳定、几何图形相似或预测准确，就证明系统拥有与人相同的主观感受。本章因此把原有波态、流态、固态、观察算子和相位等物理图景逐一落实为可观测的状态变量、更新接口与实验边界。论证从工作记忆、长期记忆和回忆机制开始，继而区分在线时序、回顾时长和跨模态绑定，最后讨论感受质报告、观察者模型及意识边界。HSF--HD在此提供现象学相关行为的一般状态动力学；AgentOS只是其中一种工程实例，而不是体验本身的定义。

\section{记忆机制：活动维持、证据积累与结构提交}
\label{sec:section_9311}
我们先把记忆定义为系统在时间间隔后仍能使过去事件影响当前任务的能力，而不是存放在单一容器里的实体。离散时刻 $k$ 的记忆状态写为 $m_k=(w_k,e_k,s_k,p_k)$：$w_k\in\mathbb R^{d_w}$ 是当前任务使用的工作状态，$e_k$ 是带时间、来源和对象身份的事件记录，$s_k$ 是跨任务保持的语义结构，$p_k$ 是程序、技能或策略参数。四者的尺度、可写权限和失效方式不同。一次强激活可以扩大 $w_k$，却不自动产生可靠的 $e_k$ 或永久改变 $s_k$；结构长期存在也不意味着维持它完全没有物理成本，因为载体刷新、校验和环境保存仍消耗资源。

工作记忆是有容量与保持时间的受控活动状态。给定输入编码 $z_k$、任务规范 $\Theta_k$、门控 $g_k\in[0,1]^{d_w}$ 和衰减矩阵 $A_k$，一种离散更新为
\[
w_{k+1}=\Pi_{\mathcal W}\!\left(A_k w_k+g_k\odot B_k z_k+(1-g_k)\odot R_k(w_k,\Theta_k)+\xi_k\right),
\]
其中 $\Pi_{\mathcal W}$ 把状态投影到容量与安全可行集，$R_k$ 是复述或重入更新，$\xi_k$ 是噪声。若连续近似满足 $\dot w=-\Lambda w+u(t)$，离散实现仍需使步长与 $\Lambda$ 的谱相容；不能只凭连续衰减公式宣称运行稳定。注意力提高的是选定内容的更新和读出权重，不是抽象“负熵泵”。工作记忆失败可表现为容量饱和、干扰、门控错误或来源丢失，需要分别测试，而不能都归因于能量不足。

从瞬时活动到长期保留，需要证据积累和结构提交。系统为候选记忆建立条目 $c=(i,v,\rho,\tau,q)$，分别记录对象身份、内容、来源、有效时间和置信度；慢速学习器在事件窗口 $E_{a:b}$ 上生成结构补丁 $\Delta s$，验证器检查重复证据、来源独立性、与既有知识的冲突、任务收益和遗忘影响。只有验证通过，才执行
\[
s_{n+1}=\operatorname{Commit}\!\left(s_n,\Delta s,V_s\right),
\]
其中 $n$ 是比工作记忆时钟更慢的提交索引。这个过程保留了“反复或强烈事件更可能留下痕迹”的经验事实，但不把强度直接当真值。创伤性事件可能极易被提取，却在时间、来源或细节上失真；低频但经可靠记录验证的事实也可以稳定保存。所谓刻蚀阈值因而由证据、任务和风险共同决定，不是统一介质屈服强度。

语义记忆保存可跨情境复用的关系、属性和规则。局部结构网络适合维护对象身份、类型关系、出处与版本，全局分布表示适合维护相似性、典型性和不确定候选；两者通过任务键、绑定索引和读出接口联合。学习“桥梁通常连接河岸”可以改变关系预测，但不能把某座具体桥当前开放的事件记录抹成无时间条件的规则。相反，把所有经历逐条保存在事件库中也不能自动形成可迁移概念。语义压缩需报告它保留了哪些查询、丢失了哪些例外以及在新域上的误差。

情景记忆不是对完整世界线的无损全息压缩，而是对一次事件序列的选择性记录。我们将事件包写为 $e_j=(\tau_j,o_j,a_j,y_j,i_j,\rho_j)$，包括时间戳、观测、动作、环境回执、身份绑定与来源元数据。压缩器 $C$ 生成摘要 $\tilde e=C(e_{1:m};\Theta)$，解码或回放器只能在给定任务下近似恢复所需成分。若两个序列产生相同摘要，它们对未保留问题就不可区分；因此回忆中的连贯叙事可能来自当前模型补全，而不是过去被逐帧重放。验证应把原始记录、后来推断和当前重构标为不同证据层。

程序记忆表现为策略在重复任务中的可用性提高。熟练动作可能减少在线搜索和工作状态占用，但“更省计算”不等于零能耗，也不保证迁移到改变后的环境。我们用成功率、响应时延、干预次数和新条件下的退化衡量熟练，而不是把它描述为沿测地线无损滑行。若固定动作序列在旧场景快速执行、遇到新障碍却无法停止，它是脆弱自动化而非更高级记忆。

遗忘是记忆系统的组成部分。工作状态的自然衰减、事件的有效期、结构冲突后的降权、容量驱动的压缩和出于隐私要求的删除具有不同机制。删除索引不必然删除外部副本，降低检索权重也不等于结构已经消失；因此每种遗忘都要给出作用域和可验证回执。对不可逆删除，系统还要检查依赖任务与法律保留义务，避免把“优化容量”当作擅自清除责任证据的理由。

在AgentOS实例中，$h(t)$可承载工作状态摘要，SPC保存对象、事件与来源关系，外部记忆 $M_e$ 保存可审计记录，TCE和CCM分别处理检索使用与慢速提交，Boundary限制敏感写入和删除，Offloading执行外部存取。这些名称不是记忆的一般必要条件。只要实现能够区分活动、事件、结构和策略，维护身份与来源，并让写入、检索、整合和遗忘都有环境回执，就可以实例化同一模型族。

本节验收覆盖保持、写入、提取、干扰和遗忘。我们改变复述门控，检查工作状态的保持曲线；注入高强度伪事件，系统不得绕过来源验证写入永久结构；提供相互矛盾的事件，要求保留冲突与版本而非平均成一个事实；测试相似但身份不同的对象，防止语义近邻覆盖具体记录；改变环境规则，检查熟练策略能否停止和重新学习；发出删除请求，核验所有声明副本的回执。报告保持时长、容量、来源校准、结构提交错误率、检索准确率、干扰、迁移退化、删除覆盖和实际资源消耗。通过这些指标，我们才能说明过去确实以可追踪方式影响当前，而不能由“几何结晶”隐喻直接推出永久、真实或主观体验。
